\chapter{Pemodelan ASR}
\label{bab:pemodelan}

\section{Pemilihan teknik}

Kedua model memakai wav2vec 2.0 XLSR-53 dengan keluaran karakter. Pada prapelatihan wav2vec 2.0, sebagian representasi audio laten ditutupi. Model dilatih untuk membedakan target terkuantisasi yang sesuai dari kandidat lain. XLSR memperluas pendekatan ini ke audio lintas bahasa dengan representasi laten bersama \cite{baevski2020,conneau2021}. Dalam tugas ini, penyetelan menggunakan pasangan audio dan transkripsi Sunda dimulai dari model yang sudah memiliki kepala ASR. Pendekatan ini memanfaatkan representasi akustik dari tahap prapelatihan. Dengan demikian, model dapat diuji pada korpus lokal yang memiliki sekitar 128 menit audio latih.

\section{Arsitektur dan fungsi objektif}

Ekstraktor fitur tujuh lapis konvolusi mengubah sampel gelombang menjadi urutan vektor. Setiap lapis menghasilkan 512 kanal, dengan ukuran kernel $(10,3,3,3,3,2,2)$ dan langkah $(5,2,2,2,2,2,2)$. Langkah total 320 sampel menghasilkan sekitar satu bingkai setiap 20 milidetik pada 16 kHz. Setelah proyeksi fitur, enkoder Transformer dengan 24 lapis, dimensi tersembunyi 1.024, dan 16 kepala atensi per lapis membentuk representasi kontekstual. Lapisan linear menghasilkan logit per bingkai untuk setiap token CTC \cite{cahyamodel,nlpmodel,transformers}.

\begin{figure}[htbp]
  \centering
  \small
  \setlength{\fboxsep}{5pt}
  \fbox{Gelombang 16 kHz}
  \(\rightarrow\)
  \fbox{CNN}
  \(\rightarrow\)
  \fbox{Transformer}
  \(\rightarrow\)
  \fbox{Kepala CTC}
  \(\rightarrow\)
  \fbox{Dekode}
  \(\rightarrow\)
  \fbox{Teks}
  \caption{Alur inferensi wav2vec 2.0 untuk satu ujaran.}
  \label{fig:alur-asr}
\end{figure}

\textit{Connectionist Temporal Classification} (CTC) melatih model tanpa penyelarasan waktu pada tingkat karakter. Untuk audio \(x\) dan transkripsi \(y\), CTC menjumlahkan probabilitas semua jalur bingkai \(\pi\) yang dipetakan menjadi \(y\). Pemetaan \(B\) menggabungkan token identik yang berurutan, lalu menghapus token kosong (\textit{blank}). Token kosong di antara dua karakter identik mempertahankan kedua karakter pada keluaran \cite{graves2006}.
\begin{equation}
  \mathcal{L}_{\mathrm{CTC}}(x,y)
    = -\log \sum_{\pi\,:\,B(\pi)=y}
      \prod_{t=1}^{T} p(\pi_t \mid x).
\end{equation}

\clearpage
\section{Model pralatih dan kosakata}

\begin{description}
  \item[Cahya.] \kode{cahya/wav2vec2-large-xlsr-sundanese}, revisi \kode{beffe9e905ecd0787c3e87f271f8df7142f23b5b}. Model ini telah disetel untuk bahasa Sunda. Kepala CTC bawaannya memiliki 29 keluaran dan kosakatanya belum memuat \textit{é}. Sebelum penyetelan pada korpus gabungan, token \textit{é} dan satu baris baru pada kepala CTC ditambahkan. Bobot keluaran yang sudah ada disalin, sedangkan baris tambahan diinisialisasi secara acak \cite{cahyamodel,cahyaresep,repoasr}.
  \item[Indonesian NLP.] \kode{indonesian-nlp/wav2vec2-indonesian-javanese-sundanese}, revisi \kode{e5e699fa5aa5bdce999276a90a29589587b58ac9}. Model ini telah disetel menggunakan data bahasa Indonesia, Jawa, dan Sunda. Kosakata bawaannya sudah memuat \textit{é} dan kepala CTC memiliki 30 keluaran. Oleh karena itu, model ini dapat disetel langsung tanpa perluasan kosakata \cite{nlpmodel,repoasr}.
\end{description}

Kedua model hasil penyetelan memiliki 30 keluaran CTC. Token \kode{|} menandai batas kata. Token pengisi \kode{[PAD]} juga berfungsi sebagai \textit{blank} CTC. Token \kode{[UNK]} menandai karakter di luar kosakata.

\section{Konfigurasi penyetelan}

\begin{table}[htbp]
  \centering\small
  \caption{Konfigurasi pelatihan penuh kedua model.}
  \label{tab:konfigurasi}
  \begin{tabularx}{\textwidth}{@{}lY@{}}
    \toprule
    Komponen & Konfigurasi \\
    \midrule
    Data dan \textit{seed} & 1.416 ujaran latih, 442 ujaran validasi, \textit{seed} 42 \\
    Ukuran \textit{batch} & 4 per GPU, 2 GPU, akumulasi 4 langkah, \textit{batch} efektif 32 \\
    Laju pembelajaran & \(3 \times 10^{-4}\), pemanasan \pct{10}, peluruhan bobot 0 \\
    Durasi pelatihan & Maksimum 50 \textit{epoch}, penghentian dini setelah 5 evaluasi tanpa perbaikan WER \\
    Parameter dan memori & Ekstraktor konvolusi dibekukan, pelatihan FP16, \textit{gradient checkpointing} \\
    Regularisasi & \textit{Attention dropout} dan \textit{hidden dropout} 0,1, \textit{feature projection dropout} 0,0, \textit{layerdrop} 0,1, probabilitas \textit{time masking} 0,3 \\
    CTC dan pemilihan model & Reduksi \kode{mean}, \kode{ctc_zero_infinity} aktif, evaluasi dan penyimpanan per \textit{epoch}, WER validasi minimum \\
    \bottomrule
  \end{tabularx}
\end{table}

Konfigurasi pada Tabel~\ref{tab:konfigurasi} berasal dari \textit{notebook} penyetelan dan \kode{training_config.json} masing-masing model \cite{repoasr}. Enkoder Transformer dan kepala CTC tetap dilatih. Pembekuan hanya diterapkan pada ekstraktor fitur konvolusi.
